\section{Design Patterns}
\subsection*{UML Class Diagram Reference for Design / Implementation Phase}
\begin{figure}[h]
    \centering
    \includegraphics[width=0.9\textwidth]{Design Pattern Images/UML Class Diagram Reference 1.jpeg}
\end{figure}
\begin{figure}[h]
    \centering
    \includegraphics[width=0.9\textwidth]{Design Pattern Images/UML Class Diagram Reference 2.jpeg}
\end{figure}
\begin{figure}[h]
    \centering
    \includegraphics[width=0.9\textwidth]{Design Pattern Images/UML Class Diagram Reference 3.jpeg}
\end{figure}
\begin{figure}[h]
    \centering
    \includegraphics[width=0.9\textwidth]{Design Pattern Images/UML Class Diagram Reference 4.jpeg}
\end{figure}
\clearpage
\subsection{Part I: Behavioural Design Patterns: Observer Pattern}
\subsubsection{Observer Pattern: Motivation, Example}
\begin{figure}[h]
    \centering
    \includegraphics[width=0.9\textwidth]{Design Pattern Images/Observer Pattern: Motivation.jpeg}
\end{figure}
\subsubsection{Consequence of OO Design}
\begin{defBox}
    OOAD (Object-Oriented Analysis und Design) divides a problem into manageable \red{subtasks}
\end{defBox}
\begin{defBox}
    $\Rightarrow$ Objects are responsible for \red{subtasks}\\
    They \blue{collaborate} to accomplish complex tasks
\end{defBox}
\fatsf{Advantages:}
\begin{itemize}
    \item Each object easy to implement, understand, maintain, reuse
    \item \blue{Divide-and-conquer} approach to design possible
    \item Flexible combinations from different collaborations possible
    \item \blue{Versatility}
\end{itemize}
\fatsf{Disadvantages:}
\begin{itemize}
    \item Behavior is \red{distributed} across multiple objects
    \item \blue{Lack of clarity}(\enquote{Unübersichtlichkeit}) of \red{overall} design
    \item Any state change of one object may affect many others
    \item \blue{Interdependence}
\end{itemize}
\subsubsection{Desirable: Communication with Low Coupling}
\begin{defBox}[title=Change Propagation]
    \begin{itemize}
        \item \red{Change propagation} (of object states) can be \red{hard wired}, but this couples objects together and diminishes their flexibility and potential for reuse
        \item A flexible way is neded to allow objects to \red{inform each other about changes without strongly coupling them}
    \end{itemize}
\end{defBox}
\begin{defBox}[title=Representative Scenario]
    Separation of GUI from underlying data, such that classes defining application data and presentations can be \red{reused independently}
\end{defBox}
\begin{defBox}[title=Task]
    \red{Decouple} a data model (\blue{\enquote{subject}}) from parties (\blue{\enquote{observers}}) interested in changes of its internal state
\end{defBox}
\begin{defBox}[title=Requirements]
    \begin{itemize}
        \item \red{Subject} should not need to know details about its \red{observers}
        \item Identity and number of observers \red{not predetermined} or fixed
        \item \red{New observer} kinds can be added to the system \red{dynamically}
        \item \red{Polling} is an inappropritate (inefficient) communication pattern
    \end{itemize}
\end{defBox}
\begin{defBox}
    Fundamental problem: one-to-many communication relation among objects
\end{defBox}
\subsubsection{Observer Pattern: Pattern Structure}
\begin{figure}[h]
    \centering
    \includegraphics[width=0.9\textwidth]{Design Pattern Images/ObserverPattern:Pattern Structure.jpeg}
\end{figure}
\begin{defBox}
    Methods getState(), modifyState() of ConreteSubject are placeholders for implementing methods of a concrete subject that get or modify its state
\end{defBox}
\begin{infoBox}[title=Wie funktioniert das Observer Pattern?]
    \begin{enumerate}
        \item \fatsf{Subject (Subjekt)}\\
        Das Subject ist das Objekt, das überwacht wird. Es verwaltet eine Liste von Beobachtern und stellt Methoden bereit, um:
        \begin{itemize}
            \item Beobachter hinzufügen (\textsf{attach} oder \textsf{subscribe})
            \item Beobachter zu entfernen (\textsf{detach} oder \textsf{unsubscribe})
            \item Beobachter über Änderungen zu benachrichtigen (\textsf{notifyObservers})
        \end{itemize}
        \item \fatsf{Observer (Beobachter)}\\
        Ein Observer ist ein Objekt, das am Zustand des Subjekts interessiert ist. Es definiert eine Methode (\textsf{update}), die aufgerufen wird, wenn das Subjekt seinen Zustand ändert
        \item \fatsf{ConcreteSubject (Konkretes Subjekt)}\\
        Das konkrete Subjekt implementiert die Schnittstelle oder Basisklasse des Subjects. Es enthält den eigentlichen Zustand und die Logik zur Verwaltung der Beobachter
        \item \fatsf{ConcreteObserver (Konkreter Beobachter)}\\
        Der konkrete Beobachter imlementiert die Observer-Schnittstelle. ER definiert, wie er auf Zustandsänderungen des Subjekts reagiert
    \end{enumerate}
\end{infoBox}
\begin{infoBox}[title=Ablauf des Observer Patterns]
    \begin{enumerate}
        \item Der Beobachter (Observer) registriert sich beim Subjekt (Subject), um Benachrichtigungen zu erhalten
        \item Das Subjekt ändert seinen Zustand und ruft die Methode \textsf{notifyObservers()} auf
        \item Das Subjekt informiert alle registrierten Beobachter, indem es deren \textsf{update}-Methode aufruft
        \item Die Beobachter reagieren auf die Benachrichtigung und passen sich entsprechend an
    \end{enumerate}
\end{infoBox}
\subsubsection{Observer Pattern: Example}
\begin{figure}[h]
    \centering
    \includegraphics[width=0.9\textwidth]{Design Pattern Images/Observer Pattern: Example.jpeg}
\end{figure}
\subsubsection{Observer Pattern: Adding and Removing Observers}
\begin{codeBlock}{}
    public abstract class Subject{
        private CopyOnWriteArrayList<Observer> observers = new CopyOnWriteArrayList<>();
        public void attach(Observer obs){
            if (obs == null){
                throw new IllegalArgumentException();
            }
            observers.addIfAbsent(obs);
        }
        public void detach(Observer obs){
            if (obs == null){
                throw new IllegalArgumentException();
            }
            observers.remove(obs);
        }
        public void notifyObservers(){...}
    }
\end{codeBlock}
\begin{defBox}[title=Remarks]
    \begin{itemize}
        \item Observers are added and removed from a list of observers
        \item Attention: In amulti-threaded environment additional care is necessary, for instance, synchronization and/or use of thread-safe data structures (here: CopyOnWriteArrayList)
    \end{itemize}
\end{defBox}
\subsubsection{Observer Pattern: Change Notification}
\begin{figure}[h]
    \centering
    \includegraphics[width=0.9\textwidth]{Design Pattern Images/Observer Pattern: Change Notification.jpeg}
\end{figure}
\subsubsection{Abstract Coupling}
\begin{defBox}
    Observer pattern realizes \red{abstract coupling} between subject and observers
\end{defBox}
\begin{defBox}[title=Abstract Coupling]
    \begin{itemize}
        \item Implementation of \blue{attach(), detach(), notify()} operations need to know only observer \red{interface}
        \begin{itemize}
            \item No coupling to observer \red{implementation}
        \end{itemize}
        \item Observers can be added and removed \red{dynamically}
    \end{itemize}
\end{defBox}
\subsubsection{Observer Patterns: Consequences}
\begin{defBox}[title=Advantages]
    \begin{itemize}
        \item \red{Abstract coupling} between \blue{Subject} and \blue{Observer}
        \item Support for \red{broadcast communication} \enquote{ausstrahlen} (as opposed to polling):
        \begin{itemize}
            \item \blue{notifyObservers()} does not specify its receivers
            \item Sender does not know the (concrete) type of the receiver
        \end{itemize}
    \end{itemize}
\end{defBox}
\begin{defBox}[title=Disadvantages and Limitations (control over updates)]
    \begin{itemize}
        \item Danger of \red{update cascades} to observers and their dependent objects
        \item Update sent to \red{all observers}, but not \red{all} might be interested in \red{any} change
        \item \red{No change details}\\
        $\Rightarrow$ Observers need to find out what has changed
        \item \red{Uniform interface} for all observer updates:\\
        $\Rightarrow$ Subject cannot send optional parameters to observers
    \end{itemize}
\end{defBox}
\subsubsection{Model-View-Controller by Example}
\begin{figure}[h]
    \centering
    \includegraphics[width=0.9\textwidth]{Design Pattern Images/Model-View-Controller by Example.jpeg}
\end{figure}
\subsubsection{Concrete Subject Implementation}
\begin{codeBlock}{}
    public class CarInventory extends Subject{
        private List<Car> carPool = new ArrayList<>();
        ...
        public void addCar(Car c){
            if (!carPool.contains(c)){
                carPool.add(c);
            }
            notifyObservers();
        }
        public void removeCar(Car c){
            carPool.remove(c);
            notifyObservers();
        }
    }
\end{codeBlock}
\begin{itemize}
    \item \blue{CarInventory} is the ConcreteSubject
    \item Methods addCar(), removeCar() change ConcreteSubject's state
    \item Change notification of registered observers via notifyObserver()
    \item Update notification takes from o.update(this)
\end{itemize}
\subsubsection{Observer Update Example}
\begin{codeBlock}{}
    public class CarInventoryView extends JFrame implements Observer{
        ...
        public CarInventoryView(CarInventory inventory){
            ...
            inventory.attach(this);
            fillTable(inventory);
            ...
        }
        @Override
        public void update(Subject subject){
            fillTable((CarInventory) subject);
        }
        private void fillTable(Carinventory inventory){
            ...
            inventory.getCars().forEach(c -> model.addRow(...));
        }
        ...
    }
\end{codeBlock}
\blue{CarInventoryView}
\begin{itemize}
    \item is a ConcreteObserver
    \item implements Observer interface
    \begin{itemize}
        \item update(Subject) called by notifyObserver() of concrete Subject CarInventory, knows \red{inventory}
        \item update() causes full table refresh
    \end{itemize}
\end{itemize}
\subsubsection{Implementation Alternatives I: Who Informs Observers of State Change?}
\begin{enumerate}
    \item State changing methods notify observers (as seen)
    \begin{codeBlock}{}
        class CarInventory extends Subject{
            List<Car> pool = new LinkedList<Car>();
            public void removeCar(Car c){
                pool.remove(c);
                notifyObservers();
            }
            ...
        }
    \end{codeBlock}
    \begin{infoBox}
        Beobachter werden sofort bei jeder Änderung benachrichtigt. Dies ist einfach zu implementieren, kann jedoch ineffizient sein, wenn viele Änderungen in kurzer Zeit stattfinden.
    \end{infoBox}
    \item Client triggers update
        \begin{codeBlock}{}
            class Subject {
                private List<Observer> observers = new ArrayList<>();
                public void addObserver(Observer observer) {
                    observers.add(observer);
                }
                public void notifyObservers() {
                    for (Observer observer : observers) {
                        observer.update();
                    }
                }
            }
            interface Observer {
                void update();
            }
            class CarInventory extends Subject {
                private List<Car> pool = new LinkedList<>();
                public void removeCar(Car c) {
                    pool.remove(c); // Keine automatische Benachrichtigung
                }
                public void addCar(Car c) {
                    pool.add(c); // Keine automatische Benachrichtigung
                }
            }
            class Car {}
            class InventoryObserver implements Observer {
                @Override
                public void update() {
                    System.out.println("Inventory updated!");
                }
            }
            // Hauptprogramm
            public class Main {
                public static void main(String[] args) {
                    CarInventory inventory = new CarInventory();
                    InventoryObserver observer = new InventoryObserver();
                    inventory.addObserver(observer);
            
                    Car car1 = new Car();
                    Car car2 = new Car();
            
                    // Der Client führt mehrere Änderungen aus
                    inventory.addCar(car1);
                    inventory.addCar(car2);
                    inventory.removeCar(car1);

                    // Der Client löst die Benachrichtigung aus
                    inventory.notifyObservers();
                }
            }
        \end{codeBlock}
        Die Methode notifyObservers() wird vom Client nach Abschluss der Änderungen aufgerufen. Dadurch wird nur eine einzige Benachrichtigung gesendet, auch wenn mehrere Änderungen durchgeführt werden
\end{enumerate}
\begin{defBox}
    \red{Advantage (of 1.)}\\
    \begin{itemize}
        \item Close to actual subject state change
        \item Observers will be reliably notified
    \end{itemize}
    \red{Disadvantage (of 1.)}
    \begin{itemize}
        \item If several changes performed at once, then each triggers notification: Causing, for example, many screen refreshes
    \end{itemize}
\end{defBox}
\subsubsection{Implementation Alternatives II: Passing Update Information}
\begin{enumerate}
    \item \red{Pull Mode} (as seen)
    \begin{codeBlock}{}
        public class CarInventoryView implements Observer{
            @Override
            public void update(Subject subject){
                fillTable((CarInventory) subject);
            }
        }
    \end{codeBlock}
    \begin{itemize}
        \item \blue{Concrete observer queries}
        \item \blue{concrete subject for change}
        \item Often too coarse gained, anything might be changed (Here: full table content refresh)
    \end{itemize}
    \begin{infoBox}
        Im Pull-Modus ruft der Beobachter (Observer) die Informationen über Änderungen direkt beim Subjekt (Subject) ab, wenn er über eine Änderung benachrichtigt wird. Das Sujekt gibt keine spezifischen Details zur Änderung weiter.\\
        \fatsf{Eigenschaften:}
        \begin{itemize}
            \item Der Beobachter fragt den vollständigen Zustand des Subjekts ab (z.B. das gesamte Auto-Inventar)
            \item Geeignet für einfache Szenarien, aber möglicherweise ineffizient, da der gesamte Zustand unabhängig von der Art der Änderung neu geladen wird
            \item \fatsf{Nachteil:} Änderungen sind oft zu grobgranular, was unnötige Operationen (z.B. vollständige Tabellenaktualisierung) verursachen kann
        \end{itemize}
    \end{infoBox}
    \item \red{Push Mode}
    \begin{codeBlock}{}
        public class CarInventoryView implements Observer{
            @Override
            public void update(Subject subject, Change e){
                if (e.getKind() == Change.CarDeletion){
                    deleteRowForCar(e.getCar());
                } else {
                    ...
                }
            }
        }
    \end{codeBlock}
    \begin{itemize}
        \item Change aspect \blue{parameter} of update()
        \item Makes efficient view update possible
    \end{itemize}
    \begin{infoBox}
        Im Push-Modus sendet das Subjekt spezifische Informationen über die Änderung an den Beobachter. Der Beobachter kann die Änderung effizient verarbeiten, ohne den gesamten Zustand des Subjekts abzufragen\\
        \fatsf{Eigenschaften:}
        \begin{itemize}
            \item Das Subjekt sendet detaillierte Informationen zur Änderung (z.B. Art der Änderung und betroffene Daten)
            \item Ermöglicht dem Beobachter, die Änderung effizient zu verarbeiten (z.B. nur eine Zeile aktualisieren oder entfernen)
            \item \fatsf{Vorteil:} Präzise und effizient, besonders bei großen oder komplexen Daten
        \end{itemize}
    \end{infoBox}
\end{enumerate}
\subsubsection{Implementation Alternatives III: Concrete Subject with Subclass Hierarchy}
\begin{codeBlock}{}
    class ComplexSubject extends Subject{
        Object o = new Object();
        public void trickyChange(){
            o = new Object(); // not so tricky ...
            notifyObservers();
        }
    }
    class MoreCmplexSubject extends ComplexSubject{
        Object anotherO = ...;
        @Override
        public void trickyChange(){
            super.trickyChange(); //causes notification
            another() = ...;
            notifyObservers(); //another notification
        }
    }
\end{codeBlock}
\blue{Problem:}\\
Invocation of trickyChange() on ComplexSubject causes \red{premature notification}, because of call to super.trickyChange() in super class
\begin{codeBlock}{}
    class ComplexSubject extends Subject{
        Object o = new Object();
        public void trickyChange(){
            doTrickyChange();
            notifyObservers();
        }
        protected void doTrickyChange(){
            o = new Object();
        }
    }
    class MoreCmplexSubject extends ComplexSubject{
        Object anotherO = ...;
        @Override
        public void trickyChange(){
            super.doTrickyChange(); //avoid notification
            another() = ...;
        }
    }
\end{codeBlock}
\blue{Solution:}\\
Separate logic for model change and notification invocation into different methods (use \red{Template Method Pattern})
\subsubsection{Implementation Alternatives IV: Specifying Modifications of Interest}
\blue{Specify relevant aspect of model when registering observer}
\begin{itemize}
    \item \blue{addObserver(Observer, Aspect):} 
    Aspect encodes the kind of events of interest for \red{that} Observer
    \item In case of a state change, subject calls \blue{triggerUpdatefor(anAspect)} on itself, which informs all observers that are registered for anAspect
\end{itemize}
\blue{Different observer interfaces for different aspects} (common JAVA use case)
\begin{itemize}
    \item Different \red{observer interfaces} are available\\
    For example: \blue{ActionListener, MouseListener}
    \item Subject \red{informs} listeners on changes they are interested in\\
    For example: Mouse click on a button
    \item Subject provides \red{registration methods} for supported listener interfaces \\
    For example: \blue{addMouseListener(MouseListener)}
    \item Observers implement \red{listener interfaces} (often inner classes of view)
\end{itemize}
\subsection{Part II: Creational Pastterns: Factory Method}
\subsubsection{Factory Method: Motivation, Example}
Assume we develop a \red{framework} for appliccations that can present \red{multi-format documents} to the user\\
Need to support a wide variety of application software:
\begin{itemize}
    \item Text editor 
    \item Word processor
    \item Vector drawing 
    \item document Viewer, ...
\end{itemize}
\begin{defBox}
    The framework must be able to manage \red{different} kinds of documents with \red{common}
\end{defBox} 
\clearpage
\subsubsection{Common Functionality for Document Management}
\begin{figure}[h]
    \centering
    \includegraphics[width=0.9\textwidth]{Design Pattern Images/Common Functionality for Document Management.jpeg}
\end{figure}
\subsubsection{Factory Method: Intent, Example}
\begin{defBox}
    Declare \red{interface for object creation}, let subclasses decide class to instantiate: Factory Method lets an application class \red{defer instantioation} to subclasses
\end{defBox}
\begin{infoBox}[title=Intent (Absicht)]
    Das \fatsf{Factory Method Pattern} ermöglicht einer Klasse, die Instanziierung ihrer Objekte an Unterklassen zu delegieren. Das Hauptziel ist, eine \fatsf{Schnittstelle für die Objekterstellung} bereitzustellen, ohne dass die konkrete Klasse, die instanziiert wird bekannt sein muss. Dadurch wird der Code flexibler und leichter erweiterbar
\end{infoBox}
\begin{codeBlock}{}
    public abstract class Document{
        public abstract void open();
        public abstract void close();
    }
    public abstract class Application{
        private List<Document> docs = new ArrayList<>();
        public void newDocument(){ //manage documents
            Document doc = createDocument();
            docs.add(doc);
            doc.open();
        }
        //factory method delegates creation to subclass
        public abstract Document createDocument();
    }
\end{codeBlock}
\blue{Concrete subclasses}
\begin{codeBlock}{}
    public class TextDocument extends Document{
        //implementation of the abstract methods
    }
    public class MyApplication extends Application{
        public Document createDocument(){
            return new TextDocument();
        }
    }
\end{codeBlock}
\subsubsection{Factory Method: UML Class Diagram of Example}
\begin{figure}[h]
    \centering
    \includegraphics[width=0.9\textwidth]{Design Pattern Images/UML Class Diagram of Example.jpeg}
\end{figure}
\subsubsection{Factory Method: General Pattern Structure}
\begin{figure}[h]
    \centering
    \includegraphics[width=0.9\textwidth]{Design Pattern Images/General Pattern Structure.jpeg}
\end{figure}
\subsubsection{Factory Method: Consequences and Variants}
\blue{Consequences}
\begin{itemize}
    \item The client application code only knows the \red{Product} interface $\Rightarrow$ it works for \red{any} \blue{ConcreteProduct} class
    \item Product provides a \red{hook} for subclasses $\Rightarrow$ can provide an etended version of an object via hook
\end{itemize}
\blue{Variants}
\begin{enumerate}
    \item \red{Creator} is abstract (as seen)
    \item \red{Creator} is concrete and provides a reasonable default implementation
\end{enumerate}
\subsection{Part III: Creational Patterns: Abstract Factory}
\subsubsection{Abstract Factory: Motivation, Example}
\begin{defBox}
    How to create a family of \red{related classes} that imlement a (set of) \red{common interface}(s)?
\end{defBox}
\begin{defBox}[title=Example(Support of different database formats)]
    Requirements:
    \begin{itemize}
        \item An application must support different database formats: \red{Ability to select database format at startup time}
        \item An application must support as yet \red{unknown} database formats: \red{Make the implementation oblivious of a specific database format}
    \end{itemize}
\end{defBox}
\subsubsection{Variability with a Common Interface: Example}
\begin{figure}[h]
    \centering
    \includegraphics[width=0.9\textwidth]{Design Pattern Images/Variability with a Common Interface.jpeg}
\end{figure}
\begin{defBox}
    Interface for \blue{ResultSet} enables iteration over result of an SQL query, \red{independent of} concrete database implementation
\end{defBox}
\subsubsection{Abstract Factory: Motivation, Example}
\begin{itemize}
    \item Not sufficient for \red{full code abstraction} from concrete database Class hierarchy of 
    \begin{itemize}
        \item \blue{CallableStatements}
        \item \blue{PreparedStatements}
        \item \blue{Blobs, ...}
    \end{itemize}
    \red{for each database} are also needed
    \item Code interacting with the database can then deal with \blue{ResultSets} \red{and} \blue{SQL statements} without referring to the concrete classes, such as \blue{FirebirdResultSet}
\end{itemize}
\begin{defBox}
    \red{But:} Concrete implementation (exact subclass) must still be known when creating the statements, result sets, etc.
\end{defBox}
\begin{defBox}
    How to avoid the need to known the concrete type \red{at creation time}?
\end{defBox}
\blue{We want to avoid having to write, for example:}\\
PreparedStatement ps = new FireBirdPreparedStatement();
\begin{defBox}[title=Why?]
    \begin{itemize}
        \item Hard coding product type makes it impossible to select a database format \red{only at startup time}, with the remaining code being \red{independent}
        \item Even \red{static changes} (\enquote{find-and-replace}) are problematic, because it is easy to miss a constructor and obtain a wrong object
    \end{itemize}
\end{defBox}
\blue{Solution:}\red{Abstract} Factory
\subsubsection{Abstract Factory: General Pattern Structure as UML Class Diagram}
\begin{figure}[h]
    \centering
    \includegraphics[width=0.9\textwidth]{Design Pattern Images/Abstract Factory.jpeg}
\end{figure}
\begin{codeBlock}{}
    public AbstractProductA createProdA(){
        return new ProductA1();
    }
    public AbstractProductB createProdB(){
        return new ProductB1();
    }
\end{codeBlock}
\subsubsection{Abstract Factory: Instantiated to Example}
\begin{itemize}
    \item Abstract factory: \blue{DBConnection}
    \item Concrete factory: \blue{FirebirdConnection}
    \item Clients can create a \red{statement}, etc., without knowing the concrete DB type: Statement s = conncetion.createStmt()
    \item To interact with different databases, \blue{connection} is \red{initialized} suitably: \blue{conncetion} = DriverManager.getConnection("org.firebird.Driver");
    \item Remaining client code is \red{independent} of concrete DB type
\end{itemize}
\clearpage
\subsubsection{Abstract Factory Pattern: Product Side}
\begin{figure}[h]
    \centering
    \includegraphics[width=0.9\textwidth]{Design Pattern Images/Product Side.jpeg}
\end{figure}
\subsubsection{Abstract Factory: Consequences}
\begin{itemize}
    \item \blue{Abstracts away from concrete products}\\
    Clients unaware of concrete Products they use, even at creation time
    \item \blue{Changing product families is easy}\\
    Changing one code line can exchange a whole product family
    \item \blue{Ensures consistency among products}\\
    Product selection encapsulated in one statement, one cannot mix product types accidentally\\
    Prevents, for example, MySqlBlob from being used with a DB2 database
    \item \red{Adding unforessen kinds of products is expensive}\\
    Need to change the abstract factory and all its subclasses
    \item \red{Object creation follows non-standard pattern}\\
    Clients must know how to use \red{factory} rather than a constructor
\end{itemize}
\subsubsection{Factory Method vs. Abstract Factory}
\begin{table}[H]
    \centering
    \rowcolors{2}{\thepagecolor}{fgcolor!10!\thepagecolor}
    \renewcommand{\arraystretch}{1.2}
    \begin{tabular}{lll}
        \toprule
        & Factory Method & Abstract Factory \\
        \midrule
        Product & Single & Family\\
        Product declaration & Client & Startup \\
        Product exchange & Concrete product & Abstract product\\
        Creation & Local(Creator) & Selection of Factory\\
        \bottomrule
    \end{tabular}
\end{table}